\documentclass{ltjsbook}
\usepackage{docmute} 
\usepackage{../../myStyle} 

\begin{document} 
\chapter{導入；多変量データと心理学} 
\label{x01_introduction} 

この授業は基礎的な心理統計を修めた人向けの，応用コースになっています。基礎的な心理統計，という言葉に私が込めた意味は，
\begin{itemize}
    \item 記述統計；得られたデータの統計量を算出したり，可視化することによってデータの特徴を把握する。
    \item 推測統計；得られたデータが母集団からの標本であると考え，標本の特徴を使って母数を推測する。さらに標本の特徴から母集団について何らかの判断(意思決定)を行う。
\end{itemize}
という2点です。特に推測統計学の領域では，確率の話やさまざまな推定法，それに伴う技術などが必要ですから，これだけでも膨大な量だったのではと思います。
こうした基礎的な知識や技術を身につけると，とりあえず手元のデータを使って差があるとかないとか，どの程度効果があったのか，と言った基本的な判断はできると思います。複雑な計算式のところは機械(統計環境Rなど)がやってくれますので，出た結果だけをみて判断すれば良いのです\footnote{もちろん結果の見方が間違っていたり，拙速だったりしてはいけないなど，注意すべきところは多々あります。}。

しかし基礎的なところだけで満足していると，「はて，何がしたかったのかな」とそもそもの問題意識を忘れてしまうことにもなりかねません。とりあえず教わった(膨大な！)プロセスを経て実験結果を見てみればOKなんでしょう，というだけでは表面的な理解に止まっていると言わざるを得ません。さまざまなシーンに適用される方法論，その計算は何を意味していて，式にはどのような意味があり，式から得られるものは何をどこまで指し示しているのか，というところまで考えるのが，この講義の狙いです。数式は数式に過ぎない，というのはその通りなのですが，その数式にどのような意味があるのか，数式から得られる結果にどのような意味があるのか，を理解した上で数値(データ)を扱うようになることが目的です。数値だけに振り回される状態からの脱却，が狙いなわけです。

この授業の前半(前期)は，講義の形をとります。テーマとしては，\keyterm{因子分析}{いんしぶんせき}{Factor Analysis}を扱います。因子分析は，いわゆる質問紙調査を行った後で適用される多変量解析法の一種で，たくさんの質問項目の背後にある「因子」を見つけ出します。その因子にはXXX特性，XXX傾向といった名前がつけられ，その上で心理学的に考察することが一般的です。因子分析の結果として性格特性が得られる，などといったりするわけですから，これはもう心理学の王道中の王道，と言えるかもしれません。しかし実際は統計ソフトウェアが計算してくれて，何だかわからないまま使っている，という人も少なくありません。質問紙調査で分析して何かがわかる，というのはどういうことなのか，原理的なところから解説を始めていきます。またこの講義ではその基礎的なところ，数式レベルでしっかりと把握した上で，数値例に進みます。講義が終わる頃には，意味内容をしっかり理解したうえで因子分析を使えるようになっていることが期待されます。

この授業の後半(後期)は，実践・演習を主にした学習になります。我々が手にしたデータは，何らかのモデル・仮定のもとで生成されたものである，という考え方に立ち，\textbf{データから生成メカニズムを推測する}ことにチャレンジします。これは非常に夢のある話です。だってデータ生成メカニズムというのは，私たちの心の中の機序そのものだからです。推測に当たってはさまざまな前提・仮定が必要ですが，得られる結果は非常に含蓄に富み，さまざまな角度から心を考えさせてくれるヒントになります。線形モデルは直線的な関係でしたが，より柔軟で豊富な表現力を持つ技術へと理解を進めていきましょう。

今回は初回ですので，基礎的な心理統計で学んだことを改めて確認・復習することを中心に，今後の講義でも用いられる数学的記号の準備をします。

\section{正規線形モデルの世界}

因子分析法や(重)回帰分析では，基本的に扱う変数が複数あります。これまでの相関・回帰分析や，実験計画の中では，説明変数と被説明変数がひとつずつ，という二変数の世界でした。これが多くなったシーンは，一般に\keyterm{多変量解析}{たへんりょうかいせき}{Multivariate Analysis}と呼ばれます。変量あるいは\keyterm{変数}{へんすう}{Variables}は，ケースごとに変わる数のことを指します。変わる「数」と言っていますが，この後述べるように数字以外のものも数字として扱いますので，「ケースごとに変わるもの」の総称だと思ってください。多変量はそれが多くあるもの，データセット全体を指します。\textbf{たくさんのデータセットの中から意味ある情報を引きだすこと}，これが多変量解析の目的です。

イメージとしては，表計算ソフトのスプレッドシートに数字がずらっと並んでいる世界です。例えば性格検査の尺度の一種である，YG性格検査は被験者に120の項目について回答させます。ひとりにつき120の変数があり，これを何百，何千人に対して実施するのですから，非常に大量のデータになっているわけです。スプレッドシートにデータを入れていくときは一般に，一行(横方向)に1ケース(1オブザベーション，1個人)であるようにし，列方向(縦方向)に変数を並べるのが一般的です。数学記号では次のように表現します。

\[ x_{ij }\]

ここで$i$は個人を，$j$は変数を指します。例えば$x_{13}$で第1番目の個人(ケース)の第3番目の変数(についての値)を意味するわけです。このように一般化することで，120ある変数でも何千分ものデータでも1つの記号で表現できますね。

さて，このような大量のデータを今から分析していくわけですが，どのような方法があるでしょうか。
データ分析の領域には，さまざまなモデル，手法があります。数え方にもよりますが，何百という種類の分析法があるかもしれません。もちろん似通ったものもありますし，同じ目的に使う異なる手法もあります。これを大きく分けるなら，まず「線形モデル」と「非線形モデル」に分割できます。

\paragraph{線形モデル}線形モデルは，回帰分析や要因計画などが含まれます。関係が直線的であること，つまり変数に2乗，3乗の項が入っていないので，同じパターンで先々まで考えることができる，単純な関係です。線形というのは$y=ax+b$のグラフを書くとわかるように，直線的な関係になることからきています。変数が$x,y$だけでなく,$y = ax + bz$のように2つあったとしても，グラフでは線が面になるだけで，ある断面で見ると直線関係であることに変わりはありません。実際の多変量データではたくさんの項が含まれ，関数全体を可視化することは不可能ですが，次元が多くなっても直線関係であることに変わりはありません。一般に線形モデルは次の形で表現されます。

\[ y = a_1 x_1 + a_2 x_2 + \cdots + a_n x_n + b\]

ここで$a_m$は第$m$番目の変数$x_m$につく\keyterm{係数}{けいすう}{coefficients}であり，変数の重要性を示す数字です。足し合わせる時に，変数の重要性を変えるものなので\keyterm{重み}{おもみ}{weight}と呼ばれることもあります。この\textbf{重みつき線型結合}が線形モデルの基本です。

線形モデルの代表的な分析方法としては，回帰分析，重回帰分析，パス解析，階層回帰分析，因子分析，主成分分析，共分散分析，判別分析などが含まれます。また今あげたモデルを全て含んだ表現形式である\keyterm{構造方程式モデリング}{こうぞうほうていしきもでりんぐ}{Structural Equation Modeling}があります\footnote{これは\keyterm{共分散構造分析}{きょうぶんさんこうぞうぶんせき}{Covariance Structural Analysis}と呼ばれることもあります。}。構造方程式モデリングは総合的な表現方法で，先にあげたモデルを下位モデルとして含むものですから，これをしっかり学べば各手法をいちいち学ぶ必要がない，とも言える究極的なモデルです。本講義では第\ref{chapter12}講で触れることになります。

線形モデルの中には他にもグラフィカルモデリング\citep{宮川1997}などが含まれますが，それは本講の範囲を超えるので専門書に譲るとしておきますが\footnote{この方法は類似の名称があること，あまりメジャーな分析方法でないこともあって，専門書もすくないのですが，技術としては非常に興味深い手法です。SEMが線形モデルの王道であり線形関係を正面から捉えているのに対し，グラフィカルモデリングは関係を裏から捉える裏線形とでもいうべき手法です。もう少し言葉を足すと，SEMがどこに線形関係があるのかを探っていくのに対し，グラフィカルモデリングはどこに関係がないか，どことどこの関係が弱いか，というところを探して無意味な関係を切断していき，残った関係が考察すべきものだ，と考えるのです。}，さらに言えばこれら線形モデルのほとんどは\keyterm{正規分布}{せいきぶんぷ}{Gaussian Curve}を仮定した確率モデル群だと言えます。合わせて正規線形モデルといいます。


\paragraph{非線形モデル}
(正規)線形モデルは，変数間関係を直線的なものだと考えるのでした。しかし，世の中のことは必ずしも線形関係ではありません。むしろ線形でない関係の方が一般的でしょう。線形関係というのは$A \to B$のように，「こうすれば，こうなる」というわかりやすい関係ですが，人間の場合は特に「叩いたら，泣く」「優しい言葉をかければ，喜ばれる」といったことでも成立しないことがいっぱいあるわけです。叩かれた人が，強がって見せるためにグッと我慢するとか，優しい言葉に絆されて泣いてしまうといった人間の感情の機微は，本当に興味深く複雑な仕組みですよね。

線形モデルは，現状に当てはまらないこともありますが，わかりやすさを優先して作られたモデルです。それに対して，現状に当てはまることを目的にすると，とても線形の関係では無理です。そこで非線形な関係でもいいから，データに適したモデルはないか，と考えられているのがこの非線形モデルです。非線形だからといって，曲線である，というだけではありません。例えば条件分岐のように，枝分かれしていくような関係なども含まれます。

非線形モデルの代表的な分析方法としては，決定木，ランダム森，サポートベクターマシン，ニューラルネットワーク，ベイジアンネットワーク，アソシエーションルール，自己組織化マップなどがあります。入門書としては\citet{豊田2008}などが網羅的で良いですが，より専門的には\keyterm{機械学習}{きかいがくしゅう}{Machine Learning}などのキーワードで選書すると良いでしょう。ここでいう学習とは，データに合わせて重みを調節する方法を指し，機械が自動的に重みを調節していく様を「学習している」と表現しているわけです。巷でA.I.(人工知能)と呼ばれているものはこうした手法の総称で，データに適していることを目的にしているので，パターンがわかれば行動の予測に使えます\footnote{線形回帰モデルで予測しただけでも，知らない人にとっては人工知能=機会が計算した予測式だ，と思ってしまっているかもしれません。}。行動の予測ができれば，例えば小売業では売り上げに直結する戦略が取れるでしょうし，犯罪者のプロファイリングなど，応用的側面はいろいろ考えられます。

じゃあもう非線形モデルが最強じゃないか，と思うかもしれませんが，この系列の総合的な問題点は「機械がなぜそのような予測をしたのか説明できない」という点にあります。機械には学び方を教えていますが，どう学んだかは機械次第なのです。理屈はわからなくても正解が出せる，というのは実用的にはいいのですが，科学的な研究の場合は少し困ります。機械が勝手に人の心を「うんうん，わかりますよ」と言ったとしても，「じゃあどうわかったのか，理屈を教えてください」といっても答えられなかったり，同じアルゴリズムでも違う機械が学習すると違う重み係数になったりして一般的な理屈が出てこないということがあります。心理学は実践的な側面もありますが，究極的には人間行動の理論を探しているので，そういう意味では非線形モデルは向いていないのです。

\paragraph{モデルの展開と全体像}

多変量解析の分析方法を，線形モデルと非線形モデルに分割して説明してきました。一言でいうなら，線形モデルは「当てはまらないこともあるけど，理屈で説明がつくモデル」であり，いわば理屈が先，現象が後です。非線形モデルは「理屈はわからないけど，データには当てはまるモデル」であり，いわば現象が先，理屈が後なのです。

どちらもデータとの当てはまりを基準にしてはいますが，その背後の考え方が違っているわけですね。(正規)線形モデルも非線形モデルも，制約を増やしたり減らしたりしながら限界とされている点を克服するべく，日々モデルの改良がなされています。

たとえば線形モデルの世界でも，正規分布以外の確率分布を仮定できます。それらは\keyterm{一般化線形モデル}{いっぱんかせんけいもでる}{Generalized Liner Model}と呼ばれ，さまざまな確率分布を仮定した上で，線形関係を見出す手法です。心理学の研究対象としているデータは，目に見えない心的状態を対象とすることが多いので，正規分布を仮定することが一般的でした。もちろん数学的に，正規分布を仮定するモデルの方が単純な形になるので，そちらの発展が先に進んだという事情もあります。また，正規分布以外の形をするデータがなかったわけではありません。所得のデータは対数正規分布のようになりますし，比率のデータは0から1までの範囲にしか値を取らないので平均が0.5からズレれば左右対称の形にはなりません。友人の数を数える時のように，0，1，2...と正の整数しか取らない離散的なデータというのもあります。しかし分析モデルが正規分布を仮定したものしかなかったので，これまではデータを正規分布の形になるように変換して分析する，ということが行われてきました。今は一般化線形モデルを使って，データの形式にあった分析をすることが基本です\footnote{古い論文を読むと，正解率のような比率のデータに対して分散分析を行ったりしていましたが，このような理由から今では推奨されません。}。

このように，正規分布の線形モデルが非常に多くあるのですが，その制約を外す方向で統計モデルが展開してきているわけです。制約の外し方としては，ここで挙げた「正規分布以外の確率モデルを使う」ということもありますし，分布の歪度・尖度といったより高次の\keyterm{積率}{せきりつ}{moment}を使う方法があります\footnote{積率について，詳しくはこの後の\ref{moment}節で。}。

また，\keyterm{数理モデリング}{すうりもでりんぐ}{Mathematical Modeling}というアプローチもあります。これは線形の仮定を外し，理屈の通る変数関係を数式で表現してデータにフィットさせるというアプローチです。当然非線形な関係になりますし，正規分布以外の確率分布も使います。この時，さまざまな確率分布が入れ子になったモデルになるので，推定方法としてはベイズ法を使うことになります。ベイズ法によるモデリングアプローチは最近の計算機技術革新によってとても身近なものになりました。この授業でも第\ref{chapter16}講から扱うことになります。

さて，多変量データを扱う世界の全体像がわかったところで，まずは線形モデルの世界から少しずつ進んでいくことにしましょう。そこで，統計モデルの種類にも大きく関わる尺度水準や，記述統計量について，改めて説明を加えておきたいと思います。

\section{尺度の四水準}
心理統計のテキストは何を開いても，まず\citet{stevens1946}による\keyterm{尺度水準}{しゃくどすいじゅん}{level of measurement}についての言及があります。何を測定した数値であるかとは別に，その数値にどのような算術処理を施すことができるかによって，数字を4つのレベルに分けるのでした。

\paragraph{名義尺度水準} \keyterm{名義尺度水準}{めいぎしゃくどすいじゅん}{nominal scale}は数字と対象が1対1で対応していることだけが重要です。男性を1，女性を2とコード化するようなもので，この時「女性は男性の2倍である」といった数としての意味はありません。男性を0，女性を42としても本質的に変わりがないからです。このような名前だけの数字は，計算ができませんので，せいぜい1が何件あったかという度数を数えて集計するにとどまります。

とはいえ，数字が直接対象を指し示しているわけですから，もっとも意味のある数字かもしれません。

\paragraph{順序尺度水準} \keyterm{順序尺度水準}{じゅんじょしゃくどすいじゅん}{ordinal scale}は，数字が大小関係の意味を持っているものです。レースで1位2位と順番がつくと，1位のほうが2位より優れていることがわかります。人間の心理的な反応，とくに5段階や7段階で評定させる心理尺度は，この水準に相当します。選好の順序は明確でも，量的な違いがわからないからです。

\paragraph{間隔尺度水準} \keyterm{間隔尺度水準}{かんかくしゃくどすいじゅん}{interval scale}は，数字と数字の感覚が等しいことが制約として加わります。たとえば気温で10℃と20℃の差は，25℃と35℃の差に等しいと言えます。これは摂氏が氷点を0℃，沸点を100℃としたうえで百等分したという定義から明らかなことです。間隔が整っているので加法・減法の計算は可能ですが，原点が定かでないので比を考えることはできません。例えば10℃は20℃の倍の熱量を持っている，とは言えないのです。なぜでしょうか。例えば，同じエネルギー状態を別の温度体系に置き換えてみたとしましょう。新しい温度体系は，氷点が100で沸点が200だったとします。そうすると10℃は110，20℃は120に該当しますが，120は110の2倍にはなっていないからです。

\paragraph{比率尺度水準} \keyterm{比率尺度水準}{ひりつしゃくどすいじゅん}{ratio scale} は，さらに絶対0点の制約を付け加えたものです。これで原点からどれ位離れているか，を基準にして計算ができますので，乗法・除法もできることになりました。物理的な単位系はこの尺度水準にあるため，緻密な計算モデルを作ることができるのですね。

さて早足で4つの水準について説明をしてきました。これが重要なのは，尺度水準によってできる計算が変わってくる点にあります。名義尺度水準は数え上げぐらいしかできません。順序尺度水準も同様で，名義や順序といった\keyterm{質的変数}{しつてきへんすう}{categorical variables}の場合，たとえば代表値を求める時も度数を数えて最頻値を報告する，というぐらいがせいぜいなのです。

これに対して，間隔尺度水準や比率尺度水準の\keyterm{量的変数}{りょうてきへんすう}{numeric variables}では，加減乗除の計算ができますので，平均値を求めたり標準偏差を求めたり，ということができるようになります。

先ほど心理尺度は順序尺度水準でしかない，という話をしましたが，質問紙調査の研究例では尺度平均点を出したり，さらに進んだ統計手法で分析したりします。実はそれができるようになるためには，尺度につけられたカテゴリー(「非常に当てはまる」「どちらとも言えない」など)を，尺度値(「非常に当てはまる」を5，「やや当てはまる」を4とする，など)にする作業を経ており，その時に「非常に当てはまる」を5とするのはなぜか，という理屈が必要です。それについては心理尺度の作成の折に詳しく説明しますが\footnote{第\ref{chapter3}講}，少なくとも盲目的に行っているわけではないことに注意が必要です\footnote{残念ながら，この辺りの理屈を気にせずに分析する人が少なくありません。そんな人に，「数字では心がわからない」なんて言って欲しくありません。}。

尺度値の付与の仕方は色々考えられます。順序尺度水準でとられたデータであっても，その背後に連続的な心理的実態があると考えてその時の相対的な大きさをつけることもできます\footnote{因子分析法の一種，\keyterm{段階反応モデル}{だんかいはんのうもでる}{Gradad Response Model}と呼ばれる手法がこれにあたります。}。あるいは0か1かという二値データであっても，それを重みづけて合算することで連続的な値にすることもできます\footnote{これはテスト理論の考え方です。0が誤答，1が正答としてテストの点数から学力を考えるのですね。}。さらに名義的な尺度水準であっても，データ全体なかで直線的な関係が最も大きくなるように数字を与えることもできます\footnote{双対尺度法という考え方がこれにあたります。詳しくは\citet{西里静彦2010}を参照。}。

このように，尺度水準がかわると，適用できる分析モデルが変わります。例えば因子分析は間隔尺度水準以上のデータにしか適用できませんが，順序尺度水準のデータであれば因子分析ではなく，段階反応モデルのようなカテゴリカル因子分析を適用しなければなりません。名義尺度水準のデータであれば双対尺度法を適用しなければなりません。間隔尺度水準以上のデータに対して，双対尺度法を適用することは可能ですが，その逆は不可能です。例えば身長のデータとして\{ 170，175, 165 \}というのがあったときに，連続変数として平均値を求めることもできますし，名義尺度水準として各一件とカウントすることもできるのですが，男性2名と女性1名がいたので平均1.3の性別があるというのは意味をなさないからです。

データがどの水準にあるのかを見極められないと，間違えた分析をすることにもなりますので，注意してください。

\section{平均と分散}
\label{moment}
間隔尺度水準以上の数字であれば，平均値や標準偏差の計算が可能です。（算術）\keyterm{平均}{へいきん}{mean}は次の式によって計算されるのでした。

\[  \bar{x_j} = \frac{1}{N} \sum_{i=1}^N x_{ij} \]

総和の記号である$\sum$の使い方などを再確認しておいてください。また変数$j$の平均は$\bar{x_j}$と表しますが，$\sum$の記号の中では$i=1$から$N$までと，$i$だけが変化しています。 $j$は変化していません。

$\sum$の記号はこの後も所々出てきます。計算の際に次のような変形をすることがありますので確認しておいてください。

\[ \sum_{i=1}^N cx_i = (cx_1 + cx_2 + \cdots + cx_n) = c(x_1 + x_2 + \cdots + x_n) =  c\sum x_i \]
\[ \sum_{i=1}^N (x_i + y_i) = (x_1 + y_1 + x_2 + y_2 + \cdots) = \sum x_i + \sum y_i \]

続いて\keyterm{分散}{ぶんさん}{variance}の式を確認しましょう。

\[ s_x^2 = \frac{1}{N}\sum(x_i - \bar{x})^2 \]

言葉でいうなら，平均偏差($x_i - \bar{x}$)の二乗の平均です。平均偏差は各点が平均点からどれぐらい離れているかを表します。その平均をとる操作($\frac{1}{N}\sum$)なので，平均的にどれぐらい離れているかがわかるのですが，そのまま平均偏差の平均をとるとゼロになりますので\footnote{平均値が全部のデータの真ん中に位置するように撮られた指標だから当然です。}，二乗するものをその値にするのでした。

この式は次のように展開できます。

    \begin{align*}
&        s_x^2  =  \frac{1}{N}\sum(x_i - \bar{x})^2 & \mbox{定義より}\\
&                  =  \frac{1}{N}\sum(x_i - \bar{x})(x_i - \bar{x}) &\mbox{}\\
&                   =  \frac{1}{N}\sum(x_i ^ 2 - 2x_i\bar{x} + \bar{x}^2) &\mbox(後ろのカッコを展開)\\
&                   =  \frac{1}{N}\sum x_i^2 - \frac{1}{N}\sum(2x_i\bar{x}) + \frac{1}{N}\sum \bar{x}^2 & \mbox{$\sum$記号の分配}\\
&                   =  \frac{1}{N}\sum x_i^2 - 2\bar{x}\frac{1}{N}\sum x_i + \frac{1}{N} N\bar{x}^2 & \mbox{$i$が変化するところだけにつく}\\
&                   =  \frac{1}{N}\sum x_i^2 -2\bar{x}^2  + \bar{x}^2 & \mbox{} \\
&                   =  \frac{1}{N}\sum x_i^2 -\bar{x}^2
    \end{align*}

計算するときは最後の形を使うことがあります。

この式で表される分散は，変数がどの程度変化しうるかということを表す値であり，変数から引き出せる情報の上限でもあります。分散が$0$，つまり変数が全く変化しなければ，どういう時に値が大きくなってどういうときに値が小さくなるのかという違いがわからないということです。違いがないものについては，それ以上考察のしようがないか，当たり前のことを言っているに過ぎないということになります。例えば質問紙調査で「他人に激しく殴られるのが好きですか」という聞き方をすると，誰しも「全く当てはまらない」と答えると思います\footnote{中にはマゾヒズムの人がいるかもしれない，という屁理屈はここでは脇に置いておいてください。}。この項目から何かがわかるか，といわれても当たり前のことすぎて何もわからない，としか言えないでしょう。これに対して「対面している相手の手足の動きが気になりますか」というような項目であれば，気にする人もいるでしょうし，気にしない人もいるでしょうから，回答に分散が生まれます。そうすると，気になった人はどういう人なのか，気にならなかった人はどういう人なのか，という考察に進むことになるわけです。このように，調査データから意味のある考察をするためには，分散の大きな項目を作らなければならないのです。

ところで，分散の式の中には二乗の項が入っていますから，このままでは元のデータと単位が異なってしまいます。そこでこの単位を整えるために，分散の正の平方根をとったものを\keyterm{標準偏差}{ひょうじゅんんへんさ}{Standard Deviation} といって散らばりの指標に使います。本講では標準偏差の記号を$s_x$と表すことにします。

平均は中心化傾向の指標の一種で，他にも中央値，最頻値などがあります。分散や標準偏差は散らばりの指標の一種で，他に最大値，最小値，範囲，IQR，パーセンタイルなどがあります。これら記述統計量は，データの特徴を記述するため，データ分析の最初のステップで確認すべき数字です。
また，分散は平均偏差を二乗したものの平均でしたが，これを三乗したものは\keyterm{歪度}{わいど}{skewness}といいます。歪度$s_x^3$の式は次のとおりです。

\[ s_x^3 = \frac{1}{N}\sum (x_i - \bar{x})^3 \]

歪度はマイナスになると左方向に，プラスになると右方向に，分布が歪んでいることを示します。ゼロに近ければ左右対称に近いことがわかります。これも記述統計量としてデータの特徴記述に利用できるでしょう。
さらに四乗したものは，\keyterm{尖度}{せんど}{kurtiosis}と言われます。

\[ s_x^4 = \frac{1}{N}\sum (x_i - \bar{x})^4 \]

尖度がゼロであれば，分布の山の尖りぐあいが平均的で，マイナスになれば潰れた山，プラスになれば尖った山の形をした分布になることがわかります。

分散が二乗，歪度が三乗，尖度が四乗でした。これらは分布の中心からどれぐらい離れているかについての累乗でしたが，これの一乗は分布の中心そのものですので，平均値がそれに当たります。これを力学の用語を借りて\keyterm{モーメント}{もーめんと}{moment}といいます。日本語では\keytermL{積率}{せきりつ}と訳されています。重心からの距離に関する指標という意味で一般化された表現です。多変量解析のほとんどは，二次のモーメント(分散)までしか活用しませんが，3次，4次のモーメントを使うとなるとデータから得られる情報が増えることになり，より表現力を増した分析ができるようになります。

\section{共分散と相関係数}
さて，分散が1変数の変動を表現する数字であり，そこから得られる情報の上限であるという説明をしました。
実際に調査的な研究をするときは，いくつもの変数を同時に扱うことになるわけですが，そうすると変数と変数の関係を考える必要があります。

改めて分散の式を見ると，次のようになっているのでした。
\[  s_x^2 = \frac{1}{N}\sum (x_i - \bar{x})^2 = \frac{1}{N}\sum (x_i - \bar{x})(x_i - \bar{x}) \]
この式を少し変えて次のようにすることで，変数$x$と$y$の関係を考える式になります。
\[ s_{xy} = \frac{1}{N} \sum(x-\bar{x})(y-\bar{y})\]
この式で表される数字を\keyterm{共分散}{きょうぶんさん}{covariance}といいます。異なる変数ですので異なる記号で現しましたが，変数を$x_{ij}$のように個人$i$と変数$j$という形で表すならば，変数$j$と$k$の共分散を表す式は次のように書けます。添字のつき方に注意して理解してください。
\[ s_{x_{jk}} = \frac{1}{N} \sum (x_{ij} - \bar{x_j})(x_{ik - \bar{x_k}}) \]

さてこの共分散は，カッコの中がそれぞれの変数の平均偏差を現しており，それを個人ごとに集めて平均しています。あるケースにおいて，平均偏差が同じ方向，すなわちどちらも平均より上であるか，どちらも平均より下であれば，積の符号は正になるのでこの数字は増えます。逆にあるケースにおいて平均偏差が異なる方向，つまり一方は平均より上なのに他方は平均より下であるようなことがあれば，積の符号は負になりますので，この数字は減ります。これを平均するわけですから，データ全体で同じ方向にずれる傾向があるのか，そうでないのかを表す指標ということになります。

同じ方向にずれるということは，一方の動きがわかれば他方の動きが推測できます。このように，共分散の数字は変数間関係から予測，考察をするための情報量を現しているとも言えます。共分散が全くない，0であれば，ケースごとに様々な可能性があるので一般的なことが言えないわけです。

この共分散は，分散と同じで単位に依存します。例えば身長と体重の共分散を計算すると，メートルとキログラムをかけ合わさった単位になります。このように，変数ごとに単位が変わると共分散同士の比較が難しくなりますので，データの標準化を考えましょう。すなわち，どのようなデータであっても単位に捉われずに比較できるようにします。スコアの標準化は次の式で行います。

\[ z_i = \frac{x_i - \bar{x}}{s_x} \]

このように平均偏差を標準偏差で割ることで，あるデータがその変数の平均からどれぐらい離れているか，標準偏差を単位とした相対的な大きさで表すことができます。標準化されたデータは，平均が$0$，分散が$1$になります。このように，標準化されたスコアは相対的に同じサイズに整えられ，単位に依存しませんので，標準化したスコア同士は相対的に比較可能です。

この標準化スコアしたスコアで共分散を計算したものが，\keyterm{相関係数}{そうかんけいすう}{correlation coefficient}です。

\[ r_{jk} = \frac{1}{N} \sum z_{ij} z_{ik} \]

この相関係数はどんな変数であっても，$-1$から$+1$の範囲に入りますので，変数間関係の相対的な比較が可能です。つまり身長と体重の関係の強さは，身長と足のサイズの関係の強さよりも大きいとかちいさい，といった表現が可能になるわけです。

多変量解析の世界は，変数がたくさんありますから，この共分散の組み合わせもたくさんあります。分散が変数から得られる情報，共分散が変数間関係から得られる情報ですから，あるデータセットから得られるこれらの情報の数はどれぐらいになるか計算してみましょう。

変数が$M$個あったとします。変数番号が$1,2,3...M$ととすると，組み合わせは1と2, 1と3....となります。このとき$i$と$i$は分散，$i$と$j$が共分散になりますが，$s_{ij}$は$s_{ji}$と同じですから，$M \times M /2$個の共分散と$M$個の分散が，このデータセットから得られる情報の全てということになります。合計で$M \times (M-1) /2 $個になりますね。このように組み合わせまで考えると，変数の数がひとつ増えるだけでも得られる情報，分析のヒントがぐっと増えることになります。多変量解析は，こうした変数間関係の情報を使って分析を進めていくことになります。

ちなみに，変数間関係を表す方法は共分散や相関係数だけではありません。これらはあくまでも変数間の直線的な関係を表す指標であることにも注意が必要です。多変量解析全体としては，変数同士の数字のズレをたしあわせた\keyterm{距離}{きょり}{distance}や，ある変数が他の変数と同時に出現した回数をカウントした共頻度などを扱うモデルもあります。


\section{課題}
\paragraph{尺度水準}4つの尺度水準の具体例をそれぞれ挙げなさい。
\paragraph{平均と分散}平均，分散，標準偏差，標準得点の定義を式で表現しなさい。
\paragraph{共分散と相関係数}共分散の定義式を書きなさい。そのとき，スコアが標準化されていると相関係数になることを示しなさい。
\end{document} 
